\section{Mixed-exponent classifications and compact images}
\label{sec:mixed-exponents}

The power-of-three theorem does not transfer to arbitrary odd
multiples by the same open-subgroup argument. We first describe the
module left by transfer alone and a nonzero localization kernel in it.
The later arithmetic arguments in this section resolve the families
$5\cdot3^r$ and $7\cdot3^r$ of Theorem~\ref{thm:lm-power-all-derived}.

\begin{proposition}[Prime-to-three transfer]
\label{prop:mixed-transfer}
Let $n=dm$ be odd, where $d=3^r\ge3$ and $(m,3)=1$, and let
$f:U_n\to U_d$ be the coordinate-power map. Put
$N_j=\pi_1((U_j)_{\bar\Q})^{(3)}$ and $A_j=N_j^{\ab}$.
Then $f_*:N_n\to N_d$ is surjective, and there is a
$G_\Q$-equivariant decomposition
\[
A_n\simeq A_d\oplus M,\qquad
M=\ker(f_*:A_n\to A_d),\qquad
\operatorname{rank}_{\Z_3}M=n^2-d^2.
\]
For a Selmer section of the full geometric pro-$3$ group of $U_n$,
subtract the compatible global cusp section whose pushdown is
given by Theorem~\ref{thm:lm-power-all-derived}. Its abelian difference
belongs to
\[
\ker\bigl(H^1(\Q,M)\longrightarrow H^1(\Q_3,M)\bigr).
\]
No vanishing of that kernel follows from transfer.
\end{proposition}
\begin{proof}
The geometric cover has degree $m^2$. A subgroup of index prime
to three maps onto every finite $3$-group quotient of the ambient
group: its image has index dividing both that index and a power
of three. Thus $N_n\to N_d$ is onto. The geometric group $N_j$
is free pro-$3$ of rank
$2g(X_j)+3j-1=j^2+1$. Homological transfer is Galois equivariant
and satisfies $f_*f^!=m^2\operatorname{id}$. Dividing by the
unit $m^2\in\Z_3$ gives the asserted splitting and rank.

Theorem~\ref{thm:localthreeintro} identifies the center at three
with the compatible rational cusp. Choose compatible parameters
$t_d=t_n^m$. Their inertia map is multiplication by the unit $m$.
Thus the global cusp parameter downstairs lifts uniquely upstairs.
Its localization agrees with the given local parameter, by the
local injectivity of peripheral Kummer cohomology. As in the proof
of Theorem~\ref{thm:restricted3}, this injectivity follows by
filling the cusp: the boundary Tate modules have no local
invariants, and the compact Tate module has none because rational
torsion on an abelian variety over $\Q_3$ is finite.

The abelian difference has zero pushdown and zero restriction
at three. The split decomposition puts it in the displayed
kernel. Transfer cannot remove $M$, on which the pushdown is zero.
\end{proof}

\begin{proposition}[A nonzero relative localization kernel]
\label{prop:mixed-kernel}
If $m>1$ in Proposition~\ref{prop:mixed-transfer}, then
\[
\ker\bigl(H^1(\Q,M)\longrightarrow H^1(\Q_3,M)\bigr)\ne0.
\]
This is a cohomological assertion. The constructed class is not
asserted to be Selmer or to lift to a nonabelian section.
\end{proposition}
\begin{proof}
Choose a prime $q\mid m$; it is odd. Write
$P_\zeta=(0:-\zeta:1)$, and let $e_\zeta$ be its inertia
generator in the boundary Tate module of $A_n$. The vector
\[
w=\sum_{\substack{\zeta^q=1\\\zeta\ne1}}e_\zeta-(q-1)e_1
\]
defines a Galois-equivariant map $j:\Z_3(1)\to A_n$.
These cusps lie over one rational cusp of $U_d$, and their
inertia generators all map with index $m$. The coefficients
sum to zero, so $j$ takes values in $M$.

The function
\[
u=\frac{Y^q+Z^q}{(Y+Z)^q}
\]
is a unit on $U_n$. Each line $Y+\zeta Z=0$ meets $X_n$
only at $P_\zeta$, with multiplicity $n$. Consequently
\[
\operatorname{div}(u)=n\left(
\sum_{\zeta^q=1,\,\zeta\ne1}P_\zeta-(q-1)P_1\right).
\]
The geometric Kummer map of $u$, restricted to $M$, gives
$h:M\to\Z_3(1)$ with
\[
h\circ j=n\bigl((q-1)+(q-1)^2\bigr)=nq(q-1).
\]

Let $a=\log_3(-2)$ and $b=\log_3(-5)$, using the injective
logarithm on $1+3\Z_3$. Both are nonzero. The global class
\[
\kappa=b[2]-a[5]\in H^1(\Q,\Z_3(1))
\]
has completed valuation $b\ne0$ at two, so has infinite order.
At three both terms have zero valuation and their logarithms
cancel; the order-two signs have zero pro-$3$ Kummer class.
Thus $\kappa|_{\Q_3}=0$. It follows that $j_*(\kappa)$ is in
the asserted kernel, and it is nonzero because
$h_*j_*(\kappa)=nq(q-1)\kappa\ne0$.
\end{proof}
